\section{Deployment pathway}\label{app:pathway}
The pathway from these components to a programme tool has two field steps, both run with health workers on the target handheld device at NP-NCD screening camps.

\para{Run-in.} About \NPathRunin{} camp women and 60 women with confirmed cancer are scanned to fix the sweep protocol, the persistence window, the quality thresholds and the operating threshold. With \NPathRunin{} women, the threshold is certified to add fewer than 5 recalls per 100 women (one-sided 95\% upper bound) when at most \NPathRuninMaxCount{} women are recalled by the AI alone.

\para{Paired accuracy study.} About \NPathCamp{} consecutive camp women (recall rate, false recalls per 1{,}000) and \NPathCancers{} women with biopsy-confirmed cancer recruited before treatment, at least 40\% of them non-palpable or $\le$2\,cm and imaging-detected (sensitivity, with this stratum and the CBE-negative cancers reported separately). CBE and the AI sweep are done in every woman. The primary comparison is CBE + AI against CBE alone and passes only if added sensitivity exceeds 5 percentage points (lower one-sided 97.5\% bound) and added recall stays below 5 percentage points (upper bound). Its power is \NPathPowerTwelve{} if the AI adds 12 points of sensitivity at 3.5\% added recall (\NPathPowerTen{} at 10 points, \NPathPowerFifteen{} at 15). With \NPathCancers{} cancers, an observed sensitivity of 0.90 has a 95\% interval of \NPathSensCiLo--\NPathSensCiHi, and the probability of showing sensitivity above 0.80 is \NPathPowerSensEighty{} when the true value is 0.90. Every recall gets radiologist ultrasound with biopsy as indicated; all camp women are followed for 12 months, and a random 10\% of screen-negative women get radiologist ultrasound to correct for partial verification. In India, ultrasound devices are regulated under the PC-PNDT Act, and use of handheld probes at camps requires registration or approval by the state authority before the run-in.
